\subsection{Sensitivity to the coefficient budget}
\label{sec:budget-sensitivity}

The coefficient budget $\kappa$ defines the prediction class; its numerical
value is not directly comparable across feature maps. An unfavorable result
at one budget does not establish impossibility at larger budgets. This
follow-up was designed after inspecting the main results and uses new
development, calibration, and test random streams.

We fix $m=12$ and compare $\kappa\in\{1,4,16\}$ on all four main systems
with new seeds 2701--2703. Other settings are unchanged, including 128 training,
512 validation, 4,096 calibration, and 2,048 test episodes per repetition.
Each budget has 21 frozen candidates. Within each system and seed, budgets
share data streams and geometric blocks; only the coefficient constraint
changes. The study comprises 36 runs and 756 models.

For each system, seed, and budget, we select the retained KAHM prefix by mean
validation MSE over the three horizons. No budget is selected using test error.
We allocate failure probability $0.05/3$ to each budget's complete family of
risk bounds and paired comparisons. A union bound gives joint confidence at
least 0.95 across the three budgets within each system--seed repetition.
This event does not extend jointly across all twelve system--seed repetitions.
Table~\ref{tab:budget} reports test errors and all-horizon certification
counts at tolerance 0.35 for the validation-selected predictors.

\begin{table}[t]
  \centering
  \small
  \setlength{\tabcolsep}{4pt}
  % Automatically generated by experiments/run_sensitivity.py.
\begin{tabular}{llrrrrr}
\toprule
System & $\kappa$ & Features & $h=1$ & $h=4$ & $h=12$ & Passes \\
\midrule
Harmonic & 1 & 16.7 & 0.3737 & 0.4102 & 0.4696 & 0/3 \\
Harmonic & 4 & 30.0 & 0.1414 & 0.2313 & 0.3900 & 0/3 \\
Harmonic & 16 & 30.0 & 0.1392 & 0.2292 & 0.3875 & 0/3 \\
\addlinespace[2pt]
Quasiperiodic & 1 & 24.7 & 0.2510 & 0.2905 & 0.3336 & 0/3 \\
Quasiperiodic & 4 & 30.0 & 0.1294 & 0.2050 & 0.2927 & 3/3 \\
Quasiperiodic & 16 & 30.0 & 0.1294 & 0.2050 & 0.2927 & 3/3 \\
\addlinespace[2pt]
Nonlinear AR & 1 & 7.3 & 0.1059 & 0.1056 & 0.1068 & 3/3 \\
Nonlinear AR & 4 & 30.0 & 0.0670 & 0.0730 & 0.0746 & 3/3 \\
Nonlinear AR & 16 & 30.0 & 0.0670 & 0.0730 & 0.0746 & 3/3 \\
\addlinespace[2pt]
Logistic & 1 & 30.0 & 0.4713 & 0.5003 & 0.5052 & 0/3 \\
Logistic & 4 & 30.0 & 0.4384 & 0.4967 & 0.5044 & 0/3 \\
Logistic & 16 & 30.0 & 0.4384 & 0.4968 & 0.5044 & 0/3 \\
\bottomrule
\end{tabular}

  \caption{Sensitivity to the coefficient budget. Test RMSE and retained
  feature count are means over three new fitting seeds. ``Passes'' counts
  validation-selected KAHM predictors certified at tolerance 0.35
  simultaneously for $h=1,4,12$. All outcomes are retained. The confidence
  allocation covers all three budget families within each system--seed
  repetition.}
  \label{tab:budget}
\end{table}

Increasing the budget from 1 to 4 materially improves fitted KAHM prediction.
On the harmonic system, mean one-step RMSE decreases from 0.3737 to 0.1414;
on the nonlinear autoregression, it decreases from 0.1059 to 0.0670. On the
quasiperiodic system, the number of validation-selected models certified at
tolerance 0.35 increases from zero to three out of three. Thus the budget affects accuracy and certification.

Increasing the budget from 4 to 16 produces much smaller changes. Harmonic
horizon-12 RMSE changes from 0.3900 to 0.3875, and none of those KAHM models
passes the all-horizon 0.35 criterion. Quasiperiodic and nonlinear-AR results
are unchanged at the displayed precision. Logistic horizon-12 RMSE remains
approximately 0.5044, with no successful all-horizon certification. Thus increasing the budget alone to 16 leaves limitations in this
configuration, without establishing impossibility for richer geometries
or larger budgets.

The autoregressive and quadratic baselines fitted at the same numerical budgets
are retained in the released results. On the harmonic system at $\kappa=16$,
their mean horizon-12 RMSEs are 0.2377 and 0.2375, respectively, and both certify
all three repetitions at tolerance 0.35. On the nonlinear autoregression at
$\kappa=4$, KAHM has smaller horizon-12 RMSE (0.0746) than the autoregressive
(0.0800) and quadratic (0.0809) baselines, while its one-step error is larger.
These task- and horizon-specific comparisons retain the feature-dependent
capacity caveat above.

Every run retains its candidate descriptions, calibration losses, pair scores,
and test losses. An independent reconstruction audit passed for all 36 runs,
recomputing 2,268 candidate--horizon certificates and 324 paired improvement
bounds. All recorded constrained optimizations converged with reported gaps
below $10^{-7}$. The use of common random streams across budgets was also
verified. The checks support the numerical calculations and confidence allocation.
